\begin{papercard}{krippendorff2004}
\cf{Stages} \stagetag{I}. A rebuttal and a guide: it corrects misconceptions in Lombard, Snyder-Duch \& Bracken's (2002) reliability guidelines, states three conditions for a coefficient to index data reliability, compares seven agreement coefficients mathematically and numerically, and ends with four recommendations. \emph{Note:} the text layer of the author manuscript drops Greek symbols (e.g.\ ``Scott's (pi)''), so quotations avoid them.

\cf{The conservative/liberal myth} Lombard et al.\ called percent agreement ``liberal'' and chance-corrected indices ``conservative'' and advised higher cut-offs for liberal ones. Krippendorff argues this treats coefficients as points on one scale, whereas they are \emph{mathematical functions} that respond to different patterns in the data \pc{krippendorff2004}{2}{I contend that trying to understand diverse agreement coefficients by their numerical results alone}. Squaring or rooting a coefficient changes only interval sizes, not what it measures. Lombard et al.'s own fallback rule (use a chance-corrected index but accept data if percent agreement is $\ge .90$) lets unreliable data through \pc{krippendorff2004}{3}{they give no empirical justification for their choice}.

\cf{Scales} Percent agreement is anchored at only one meaningful point (100\,\%); its zero means systematic extreme disagreement, which independent coders using the same instructions should not produce. The real worst case is chance coding, where chance-corrected coefficients equal 0, so these scales are anchored at two meaningful points \pc{krippendorff2004}{2}{chance-corrected agreement coefficients, by contrast, become zero when coders‟ behavior bears no relation to the phenomena to be coded}. Percent agreement also handles only nominal data and two coders, and becomes harder with more categories.

\cf{Three conditions for an index of reliability} (1) Applied to \emph{proper reliability data}: a sample from the population whose reliability is in question, described independently by widely available coders using the same instructions. (2) Units are separately describable \emph{without presupposing correctness}; coders are interchangeable and idiosyncrasies count as disagreement: ``Reproducibility is about data making, not about coders.'' (3) Its values relate to the conditions for relying on imperfect data, with at least two designated values (perfect, and chance) \pc{krippendorff2004}{3}{agreement is what we measure; reliability is what we wish to infer from it}. Many problems arise from violating independence, giving instructions coders cannot follow, or applying them to data they do not understand \pc{krippendorff2004}{3}{Many methodological problems in testing reliability stem from violating the requirement for coders to be truly independent}.

\cf{Seven coefficients in the 2$\times$2 case} Percent agreement; Osgood/Holsti CR; Bennett's $S$ (reinvented at least five times; inflated by unused categories); Scott's $\pi$; Krippendorff's $\alpha$; Cohen's $\kappa$; Benini's $\beta$. All share the observed disagreement $b+c$ and differ in expected disagreement. $\pi$ and $\alpha$ estimate expected disagreement from the \emph{pooled} proportions in the data (a population estimate that treats coders as interchangeable) and differ only by the small-sample factor $n/(n-1)$ \pc{krippendorff2004}{5}{it would make good sense, indeed, to calculate expected disagreements from the proportions of categories in that population}. $\kappa$ instead uses $p_Aq_B+p_Bq_A$, like $\chi^2$, i.e.\ it conditions on coder preferences, so it is ``a hybrid coefficient'' \pc{krippendorff2004}{6}{I suggested elsewhere (Krippendorff, 1978) that is a hybrid coefficient}; Fleiss conceded that for interchangeable coders $\pi$/$\alpha$ is correct and $\kappa$ should be restricted to one fixed pair. Lombard et al.'s claim that $\pi$ ``assumes identical distribution'' is wrong: $\pi$ and $\alpha$ merely estimate population proportions.

\cf{Numerical demonstration (Figure 3)} Three tables with the same 46\,\% agreement: with equal margins $\pi=\kappa=\alpha=.186$; with unequal margins $\kappa=.258$ while $\pi$ and $\alpha$ stay .186. $\kappa$ ``adds disagreement about marginals to the measure as an agreement'', and it cannot reach 1.0 when margins differ (maximum .505 in the third table) while $\alpha$ registers perfect agreement given the margins \pc{krippendorff2004}{7}{just the opposite is correct}, \pc{krippendorff2004}{7}{Predictability has nothing to do with reliability}. Cohen's original criticism of Scott was mistaken, and $\kappa$ is ``just about worthless as a reliability index'' for content analysis (Brennan \& Prediger, Zwick) \pc{krippendorff2004}{8}{Its violating (2) and (3) renders just about worthless as a reliability index in content analysis}.

\cf{Rare-category example} One of Lombard et al.'s variables: 83 joint absences, 3 mismatches, no joint presence. Percent agreement is 96.5\,\% (passing their $.90$ rule), while chance-corrected coefficients are near zero (.012--.024). The reported finding (2 of 137 articles) has zero concurrences, so ``the use of percent agreement should be actively discouraged, especially as a fallback criterion''; when all coders use one category there is no evidence of reliability, since reproducibility requires evidence it could be otherwise \pc{krippendorff2004}{10}{Without a single concurrence and three mismatches, the report of finding 2 out of 137 cases is about as close to chance as one can get}.

\cf{Per-distinction reliability} Coefficients are averages over categories and let unreliable categories hide. Test every distinction that matters. Example: a three-category variable with $\alpha=.686$; 0 vs.\ 1\&2 gives .789; category 1 vs.\ rest .006 (``the unquestionable culprit''); the 1-vs-2 distinction is exactly .000, so reporting differences between those categories would be a mistake \pc{krippendorff2004}{11}{If its categories were equally unreliable, one could let the overall reliability of the variable stand}. Lumping categories can restore reliability if the lost distinction is not needed.

\cf{Multiple coders and overlapping sets} Using different coders for different but overlapping units is legitimate for $\alpha$, which is designed for variable numbers of interchangeable coders and missing data; averaging pairwise agreements (as PRAM does) would average unrelated data \pc{krippendorff2004}{11}{if reliability data are generated by multiple coders of different but overlapping sets of data}.

\cf{Four recommendations} (i) Reliability data from widely available coders, with explicit, communicable instructions, working independently, representative of the data in question; coders interchangeable, different subsamples allowed if enough overlap. (ii) Use a coefficient that measures agreement within multiple descriptions regardless of number and kind of coders; for two coders, large samples and nominal data $\pi$ is one, and $\alpha$ handles all other cases; reject association measures, $\kappa$, Cronbach's $\alpha$, Goodman--Kruskal $\lambda$, percent agreement. (iii) Choose the acceptable level by the costs of invalid conclusions: for scholarly use $\alpha\ge.800$, tentative conclusions at $\alpha\ge.667$, with confidence intervals; testing a null of chance agreement is irrelevant \pc{krippendorff2004}{12}{Testing the null-hypothesis that observed agreement deviates from chance has no bearing on reliability}. (iv) Test every distinction that matters; for a system of variables the smallest reliability is the reliability of the whole; do not average with easy clerical variables; majority resolution among three coders ``does not affect the measured reliability'' \pc{krippendorff2004}{13}{Averaging the agreement measures of several variables, especially when they include easily coded clerical ones}.

\cf{Use for this thesis} Supplies the argument for the thesis' choice of agreement statistic and what to report: chance-corrected, interchangeable coders, unit-level reliability data, per-category (one-vs-rest) values for each decision that enters a result, no averaging across variables, and CI-based thresholds. The rare-category example is directly relevant to skewed privacy codes (e.g.\ rare dark patterns): high percent agreement can hide zero joint detections, so report joint presences and consider targeted sampling. It is the author's rebuttal, so state that the $\kappa$/AC1 verdicts are his position, and weigh them against Hallgren, McHugh, Gwet and Feinstein \& Cicchetti.
\end{papercard}
